\begin{table}[htbp]\centering
\def\sym#1{\ifmmode^{#1}\else\(^{#1}\)\fi}
\caption{Political Attitudes and Engagement}
\begin{tabular}{l*{1}{ccccc}}
\toprule
                    &\multicolumn{1}{c}{(1)}&            &            &            &            \\
                    &           N&        Mean&          SD&         Min&         Max\\
\midrule
Party Identification (1-7)&        8947&        3.83&        2.34&           1&           7\\
Democratic Party Thermometer (0-100)&        8947&       46.15&       33.11&          -8&         100\\
Republican Party Thermometer (0-100)&        8947&       42.50&       33.05&          -8&         100\\
Affective Polarization&        8947&       54.37&       29.00&           0&         100\\
Political Attention (1-5)&        8947&        2.22&        0.99&           1&           5\\
\bottomrule
\multicolumn{6}{l}{\footnotesize Sample: 9,093 respondents from ANES 2020 and 2024.}\\
\multicolumn{6}{l}{\footnotesize Party ID: 1=Strong Democrat, 4=Independent, 7=Strong Republican.}\\
\multicolumn{6}{l}{\footnotesize Mean party ID (3.84) indicates slight Democratic lean in sample.}\\
\multicolumn{6}{l}{\footnotesize Thermometers: 0=coldest feelings, 100=warmest feelings toward party.}\\
\multicolumn{6}{l}{\footnotesize Affective polarization: absolute difference between party thermometers.}\\
\multicolumn{6}{l}{\footnotesize Political attention: 1=always pay attention, 5=never. Mean (2.22) = most of the time.}\\
\end{tabular}
\end{table}
